\section{Приложение Б. Каталог сценариев и эксплуатационные материалы}

\subsection{Каталог функциональных сценариев}

\begin{longtable}{p{0.08\textwidth}p{0.28\textwidth}p{0.56\textwidth}}
\toprule
ID & Сценарий & Критерий приемки \\
\midrule
\endfirsthead
\toprule
ID & Сценарий & Критерий приемки \\
\midrule
\endhead
F-01 & Регистрация агента & После старта Agent появляется в списке Master с валидным \texttt{agentId} и \texttt{baseUrl}. \\
F-02 & Heartbeat агента & Поле последнего heartbeat обновляется периодически, агент не удаляется TTL-процедурой. \\
F-03 & Публикация compose-проекта & API возвращает успешный статус, на диске агента сохраняется compose-файл проекта. \\
F-04 & Запуск проекта & После команды start сервисы переходят в состояние running, статус доступен через API. \\
F-05 & Остановка проекта & После stop сервисы удаляются/останавливаются в соответствии с настройками. \\
F-06 & Чтение логов & Запрос логов возвращает поток строк по выбранному сервису без критических ошибок. \\
F-07 & Добавление TCP-маршрута (auto) & При пустом listen выбирается первый доступный порт в диапазоне master-конфига. \\
F-08 & Добавление TCP-маршрута (fixed) & При явном listen в допустимом диапазоне маршрут сохраняется и отображается в UI. \\
F-09 & Отказ при конфликте listen-порта & API возвращает информативную ошибку, состояние реестра не изменяется. \\
F-10 & Удаление TCP-маршрута & Маршрут удаляется из JSON и cfg, в UI больше не отображается. \\
F-11 & Релоад HAProxy после изменения cfg & Изменение файла приводит к HUP-сигналу и применению новых listen без ручного вмешательства. \\
F-12 & Проверка BOM-защиты & При обнаружении BOM в cfg конфигурация переписывается в UTF-8 без BOM. \\
F-13 & Доступ к backend на хосте & Маршрут \texttt{host.docker.internal:PORT} проходит health check и пропускает TCP-трафик. \\
F-14 & Просмотр сгенерированного cfg в UI & Пользователь видит актуальный текст динамического конфига для диагностики. \\
F-15 & Sandbox Fluent preview & Код в sandbox формирует YAML или возвращает контролируемую ошибку валидации. \\
F-16 & Очистка stale agents & Фоновая задача удаляет неактивных агентов по заданному TTL. \\
F-17 & Аудит операций & Критичные действия фиксируются в журнале аудита в базе Master. \\
F-18 & Повторный подъем стека & После рестарта compose состояние маршрутов и БД консистентно восстанавливается. \\
\bottomrule
\end{longtable}

\subsection{Набор эксплуатационных рисков и мер снижения}

\begin{longtable}{p{0.3\textwidth}p{0.28\textwidth}p{0.34\textwidth}}
\toprule
Риск & Вероятное проявление & Мера снижения \\
\midrule
\endfirsthead
\toprule
Риск & Вероятное проявление & Мера снижения \\
\midrule
\endhead
Несоответствие backend host топологии & TCP health check в HAProxy в состоянии DOWN & Использование \texttt{host.docker.internal} для хостовых портов, явная документация в UI. \\
Конфликт listen-портов & Отказ добавления маршрута или bind error при старте HAProxy & Диапазон listen в env, проверка занятости и резервирования портов в реестре. \\
CRLF в shell-скриптах & Ошибки \texttt{illegal option} при запуске entrypoint & Правило \texttt{.gitattributes} и контроль формата файлов в CI. \\
BOM в динамическом cfg & Ошибка парсинга HAProxy \texttt{unknown keyword} на первой строке & Принудительная запись UTF-8 without BOM и автоисправление при загрузке. \\
Некорректный HUP/reload & Маршрут добавлен в файл, но не применен процессом & Автоматический watcher-скрипт и явные health checks. \\
Проблемы доступности БД & Ошибки при старте Master/Agent или миграциях & Health checks postgres, отложенный запуск сервисов по \texttt{depends\_on}. \\
Рост аудита без контроля retention & Деградация БД по размеру таблиц & Периодическая очистка старых записей по configurable retention. \\
Неоднозначность user input в UI & Пользователь указывает контейнерный порт вместо хостового и наоборот & Контекстные подсказки и информативные ошибки API. \\
\bottomrule
\end{longtable}

\subsection{Расширенный шаблон экспериментальной главы}

Ниже приведен текстовый шаблон, который может быть использован для расширения основной главы результатов до полного стандарта кафедральной ВКР. Каждый подраздел рекомендуется дополнять таблицами замеров, скриншотами и логами запусков.

\textbf{1. Описание аппаратно-программного стенда.}
Указать характеристики хоста (CPU, RAM, диск), версию операционной системы, версии Docker Engine, Docker Compose, .NET SDK/runtime, браузера для UI-тестов и сетевые ограничения. Важно фиксировать точные версии образов (\texttt{postgres:16-alpine}, \texttt{haproxy:3.0-alpine}, \texttt{grafana/grafana:11.3.0}, \texttt{grafana/loki:3.1.1}), поскольку различия в минорных релизах влияют на воспроизводимость.

\textbf{2. Методика проведения измерений.}
Определить повторяемые процедуры: число прогонов каждого сценария, метод фиксации времени (например, timestamp логов), критерии успешности и правила отброса выбросов. Для корректного сравнения важно разделять ``холодный'' запуск (первый старт после поднятия образов) и ``теплый'' запуск (повторный запуск без пересборки).

\textbf{3. Серия экспериментов по маршрутизации.}
Эксперименты рекомендуется разделить на: а) добавление 1 маршрута; б) добавление серии маршрутов; в) удаление и повторное добавление; г) сценарий конфликтующих listen-портов; д) сценарий недоступного backend. Для каждого эксперимента фиксировать: время отклика API, время применения конфигурации HAProxy, фактическую доступность конечного backend.

\textbf{4. Эксперименты устойчивости.}
Проверить поведение при рестарте контейнера HAProxy, Master и Agent, а также при рестарте всего compose-стека. Ключевая проверка: после восстановления маршруты должны остаться консистентными в JSON-реестре, UI и runtime-состоянии HAProxy.

\textbf{5. Эксперименты ошибок ввода.}
Проверить реакции на пустой backend host, нечисловой backend port, порт вне диапазона 1--65535, listen-порт вне опубликованного диапазона, дублирование listen-порта. Для диплома важно показать не только ``happy path'', но и инженерную зрелость обработки ошибок.

\textbf{6. Анализ результатов.}
Сформулировать выводы: какие сценарии устойчивы, где система требует операторской квалификации, какие ошибки перехватываются автоматически, а какие остаются зоной ответственности инженера эксплуатации.

\subsection{Таблица API-эндпоинтов для документирования}

{\small\raggedright
\begin{longtable}{
  >{\raggedright\arraybackslash}p{0.26\textwidth}
  >{\raggedright\arraybackslash}p{0.12\textwidth}
  >{\raggedright\arraybackslash}p{0.18\textwidth}
  >{\raggedright\arraybackslash}p{0.38\textwidth}}
\toprule
Путь & Метод & Подсистема & Назначение \\
\midrule
\endfirsthead
\toprule
Путь & Метод & Подсистема & Назначение \\
\midrule
\endhead
\path{/agents/register} & POST & Master API & Регистрация агентного узла в реестре Master. \\
\path{/agents/{id}/heartbeat} & POST & Master API & Обновление heartbeat агента и метрик состояния. \\
\path{/agents} & GET & Master API & Получение списка активных и исторических агентных узлов. \\
\path{/api/v1/haproxy/tcp-routes} & GET & UI API & Чтение списка маршрутов и сгенерированного HAProxy cfg. \\
\path{/api/v1/haproxy/tcp-routes} & POST & UI API & Добавление TCP-маршрута с валидацией listen/backend параметров. \\
\path{/api/v1/haproxy/tcp-routes/{id}} & DELETE & UI API & Удаление маршрута по идентификатору. \\
\path{/api/v1/sandbox/fluent-preview} & POST & UI API & Выполнение sandbox-просмотра Fluent-кода и генерация YAML. \\
\path{/projects} & GET & Agent API & Список compose-проектов на агентном узле. \\
\path{/projects/{name}/status} & GET & Agent API & Статус конкретного проекта и контейнеров. \\
\path{/projects/{name}/logs} & GET & Agent API & Логи сервисов compose-проекта. \\
\bottomrule
\end{longtable}
}
